% Reactive Traits Under Contract: Interaction Semantics
% for Spatial Languages with Algebraic Conflict Resolution
% Target: ECOOP 2027 (submit Feb '27)
% Backup venue: Onward! 2027 (submit Jun '27)
% Format: LIPIcs (Dagstuhl)
% Status: Draft — core sections present; camera-ready literature pass pending (2026-04-18)
% Paper ID: P3-S2 (second in Program 3: The Stalk and the Center)
% Subsystem proven: packages/core/src/parser/HoloScriptPlusParser.ts + trait system + traits tree under packages/core/src/constants/traits/
% Draft: April 2026

\documentclass[sigconf,nonacm]{acmart}

% ── Packages ─────────────────────────────────────────────────────────────────
\usepackage{amsmath}
\let\Bbbk\relax % acmart/newtxmath already defines \Bbbk; prevent amssymb redefinition clash
\usepackage{amssymb}
\usepackage{booktabs}
\usepackage{xcolor}
\usepackage{listings}
\usepackage{algorithm}
\usepackage{algpseudocode}
\usepackage{tikz}
\usetikzlibrary{arrows.meta,positioning,shapes.geometric,fit,calc,
  decorations.pathreplacing,backgrounds}
\usepackage{pgfplots}
\pgfplotsset{compat=1.18}
\usepackage{subcaption}
\usepackage{balance}
\usepackage{stmaryrd}

% ── Draft annotations ────────────────────────────────────────────────────────
\newcommand{\todo}[1]{\textcolor{red}{\textbf{[TODO: #1]}}}
\newcommand{\stub}[1]{\textcolor{gray}{\textit{[stub: #1]}}}

% Lotus petal data contract — \measuredFrom{<artifact>} marks values that come
% from a runnable artifact in the repo. Renders nothing in the PDF; grep target
% for scripts/build-lotus.mjs (research/2026-04-26_lotus-petal-data-contract.md §3+§4).
\providecommand{\measuredFrom}[1]{}

% ── Convenience macros ──────────────────────────────────────────────────────
\newcommand{\holoscript}{\textsc{HoloScript}}
\newcommand{\hsplus}{\texttt{.hsplus}}
\newcommand{\traitop}[1]{\texttt{@#1}}
\newcommand{\hashfn}{\mathcal{H}}
\newcommand{\tropplus}{\oplus}
\newcommand{\tropmult}{\otimes}
\newcommand{\provsemiring}{\mathcal{S}}
\newcommand{\traitset}{\mathcal{T}}
\newcommand{\conflictop}{\diamond}

% ── Theorem environments ────────────────────────────────────────────────────
\newtheorem{definition}{Definition}
\newtheorem{property}{Property}
\newtheorem{theorem}{Theorem}
\newtheorem{lemma}{Lemma}
\newtheorem{corollary}{Corollary}

% ════════════════════════════════════════════════════════════════════════════
\title{Reactive Traits Under Contract:\\
Interaction Semantics for Spatial Languages\\
with Algebraic Conflict Resolution}

\author{Joseph Krzywoszyja}
\affiliation{%
  \institution{Infinitus Systems}
  \country{USA}
}
\email{brianonbased@gmail.com}

% ════════════════════════════════════════════════════════════════════════════
\begin{document}

\begin{abstract}
Interaction in VR/XR is written imperatively: event handlers,
collision callbacks, input polling. When two interaction
modalities collide---a user grabs an object that is simultaneously
being thrown by physics---the outcome depends on execution order,
not on declared intent. We introduce \hsplus{}, a reactive
extension of the \holoscript{} spatial language that expresses
interactions as declarative \emph{traits}: \traitop{grabbable},
\traitop{throwable}, \traitop{hoverable}, and a large shipped library
of additional traits (see \texttt{packages/core/src/constants/traits/}).
Trait composition is an algebraic operation
under the provenance semiring~\cite{green-semiring2007}: when two
traits conflict on the same transform property, the resolution is
a deterministic, hash-verifiable semiring computation---not an
implicit execution-order dependency. We prove a conflict-resolution
theorem guaranteeing that any set of applied traits produces a
unique, contracted output regardless of evaluation order. The
order-independence guarantee is a multi-implementation agreement
generalizing the binary two-implementation agreement of the W.315
wire-format equivalent primitive: every distinct trait-evaluation
order is an independent implementation of the same semiring
state machine, and the theorem witnesses deterministic replay through
identical hash outputs across all permutations, the same-adapter
strict-digest regime; cross-vendor or cross-runtime divergence is
handled by the Route~2b per-step canonical projection pattern shared
across the program~\cite{paper-0c-twin-test-2026-05-11}. We
discuss validation against the shipped \holoscript{} trait tree
under \texttt{packages/core/src/constants/traits/}; runtime overhead
versus an uncontracted imperative baseline is $0.89$--$1.80\,\mu\text{s}$ per trait (median;
batch\,1--100, Vitest~4.1 / Node~22; 
harness: \texttt{paper-trait-semiring-overhead.test.ts},
artifact: \texttt{.bench-logs/paper-trait-semiring-overhead.json}, 2026-04-18). This is the
second paper in Program~3 of a 16-paper research
program~\cite{trustbyconstruction2026,capstone-uist2027}.
\end{abstract}

\keywords{traits, interaction semantics, conflict resolution,
  provenance semiring, VR, XR, spatial computing, reactive
  programming, declarative languages}

\maketitle

% ────────────────────────────────────────────────────────────────────────────
\section{Introduction}
\label{sec:intro}

Interaction in spatial computing is written imperatively.
Unity's \texttt{OnTriggerEnter} fires in response to physics
overlap, but its timing relative to \texttt{FixedUpdate}
depends on the component execution order---an ordering that
is globally mutable and not locally declared. Unreal's
Blueprint event graphs execute nodes in topological order
within a single graph, but the sequencing of events
\emph{across} graphs is determined by the engine's tick
scheduler, not by the developer's intent. A-Frame's component
lifecycle hooks (\texttt{init}, \texttt{tick}, \texttt{play})
execute in registration order, producing initialization bugs
when components depend on each other.

These are not implementation errors to be patched. They are
symptoms of interaction models that lack compositional
semantics. When two interaction modalities collide---a user
grabs an object that is simultaneously being thrown by
physics---the outcome depends on which handler runs first,
not on which behavior was declared with higher priority. We
propose that the atomic unit of spatial interaction is the
\emph{trait}: a declarative, typed behavioral annotation
whose composition with other traits is an algebraic
operation, not a procedural race.

\paragraph{Contributions.}
\begin{enumerate}
  \item Formal semantics of \hsplus{} traits as elements of the
    provenance semiring $\provsemiring$.
  \item Conflict-resolution theorem: for any trait set
    $\traitset = \{t_1, \ldots, t_n\}$ applied to an object,
    the composed output is unique under the semiring's
    $\tropplus$ operation, regardless of evaluation order.
  \item Reactive state propagation with contracted update:
    state changes flow through the trait graph with hash-verified
    provenance at each step.
  \item Empirical validation plan: exercised trait handlers and
    composition paths in CI; library-scale counts pinned in
    \texttt{HoloScript/.bench-logs/paper-trait-inventory.json} and
    \texttt{HoloScript/.bench-logs/paper-trait-conflict-census.json} and
    \texttt{HoloScript/.bench-logs/paper-trait-conflict-families.json}
    (see \texttt{research/benchmark-results-2026-04-18-paper-11-trait-harnesses.md}).
    Semiring-resolution overhead: median $0.89$--$1.80\,\mu\text{s}$ per trait (batches of 1--100 traits, 500\,runs, Vitest~4.1 / Node~22); JSON serialisation baseline $0.70\,\mu\text{s}$ per call (p99 $21\,\mu\text{s}$). Harness: \texttt{paper-trait-semiring-overhead.test.ts}; artifact: \texttt{.bench-logs/paper-trait-semiring-overhead.json} (2026-04-18).
  \item Comparison: imperative interaction in Unity/Unreal/A-Frame
    exhibits execution-order-dependent outcomes for the same trait
    combinations; \hsplus{} does not.
\end{enumerate}

\paragraph{Novelty.}
To our knowledge, \hsplus{} presents the first algebraic trait
system for spatial computing objects in which composition
order-independence is a \emph{provably enforced invariant}
rather than a convention or lint rule. Prior work on
entity-component systems (Unity, Unreal, A-Frame) achieves
composability through runtime dispatch order; algebraic
game-object frameworks (Helm, Arrowized FRP) enforce order-independence
for value streams but not for hash-traced provenance. We show
that equipping the trait algebra with the provenance semiring
$\provsemiring$ simultaneously guarantees order-independence
\emph{and} auditable identity, properties that are conflated in
imperative systems and absent in purely functional ones.

% ────────────────────────────────────────────────────────────────────────────
\section{Background and Related Work}
\label{sec:related}

\subsection{Imperative Interaction in Spatial Engines}
Unity components receive event callbacks
(\texttt{OnTriggerEnter}, \texttt{OnCollisionStay}) whose
execution order is determined by the Script Execution Order
settings---a global, mutable priority list. Unreal Blueprint
event graphs dispatch events through a node-based visual
script whose inter-graph sequencing depends on the tick
scheduler. A-Frame components execute lifecycle hooks in DOM
registration order. React Three Fiber's \texttt{useFrame}
callbacks execute in hook-call order. None of these systems
have formal composition semantics: the outcome of applying
two conflicting interaction behaviors depends on execution
ordering, not on declared intent.
Deterministic replay and lockstep stacks~\cite{lockstep2001} tame
cross-machine divergence but still treat interaction as ordered
callbacks rather than as a compositional algebra over conflicting writes.

\subsection{Formal Interaction Models}
Interactive Cooperative Objects (ICO, Palanque et al.) model
interaction as state machines with Petri net
semantics---formally rigorous but designed for GUI widgets,
not for spatial trait composition. Reactive input stream
frameworks (UW, 2024) provide declarative stream processing
for XR input but do not address the conflict problem when
multiple streams write to the same property. Haskell-style
typeclasses provide open, extensible dispatch but lack
provenance tracking~\cite{provone2014}. None apply semiring
algebras~\cite{green-semiring2007} to interaction conflict resolution.

\subsection{Entity-Component Systems (ECS)}
Unity DOTS/Entities, Bevy ECS, and Flecs provide
data-oriented composition through systems that iterate over
entity archetypes. System ordering graphs resolve
\emph{when} systems run (scheduling), but not \emph{what}
the composed output \emph{means} when two systems write to
the same component field. The distinction is fundamental:
ordering graphs are a scheduling mechanism; the provenance
semiring is a semantic mechanism. A correctly scheduled ECS
can still produce execution-order-dependent output if two
systems write to the same field in the same frame.

\subsection{Collaborative VR Conflict Resolution}
Multi-user VR systems resolve object-level conflicts through
locking (pessimistic concurrency), action-level restrictions
(only one user can grab at a time), or averaging strategies
(blend positions from multiple users). These are multi-user
strategies for the \emph{same} interaction. Our problem is
different: resolving conflicts between \emph{different}
interaction modalities (grab vs. throw vs. hover) on the
same object by the same user.

\subsection{Literature Survey and Positioning (camera-ready)}
\label{sec:related:survey}
Trait-oriented language design has long separated composition intent from
class inheritance. The original Traits model in Smalltalk~\cite{scharli2003traits,ducasse2006traits}
argues for conflict-explicit composition; this directly informs our
\hsplus{} requirement that conflicting writes are semantically declared,
not silently resolved by callback order. Mainstream ecosystems adopted
trait-like abstraction in distinct forms: Rust traits as bounded ad-hoc
polymorphism~\cite{rustbook2019}, Scala traits as linearized mixin
composition~\cite{odersky2004scala}, and Ruby modules as behavioral mixins~\cite{thomas2004ruby}.

Extensibility mechanisms pre-dating modern game engines also frame our
position. Emerald's distributed object model~\cite{black1987emerald} and
Self's prototype semantics~\cite{ungar1987self} show that open-world
composition is viable, but neither system attaches provenance algebra to
conflicting behavioral writes. In the module-system lineage, Standard ML
modules and higher-order module calculi~\cite{macqueen1984modules,harper1994hml}
formalize large-scale composition boundaries; we reuse that spirit for
trait families and write-set declarations in spatial code.
For operational contrast, ECS manuals and functional-reactive GUI lineage
expose deterministic scheduling and stream propagation mechanisms, but still
leave same-field write composition to user policy~\cite{bevy-ecs2024,unity-dots2024,flecs-ecs2023}.
This gap between scheduler determinism and semantic determinism also appears in
reactive-programming debugging and FRP GUI lineage~\cite{salvaneschi-reactive2018,frp-elm2012}.

On the algebraic side, provenance semirings~\cite{green-semiring2007}
and workflow provenance extensions such as ProvONE~\cite{provone2014}
establish why compositional traceability matters. Our contribution is to
carry semiring composition into interaction conflict resolution itself,
not only into offline lineage reporting. For spatial declaration prior
work, VRML/X3D-style scene description~\cite{carey1997vrml,brutzman2008x3d}
and modern declarative scene frameworks (A-Frame/R3F style APIs) support
declarative structure, but do not provide order-independent semantics for
same-property multi-trait writes.

Game engines document callback ordering (Unity script execution order,
Unreal tick groups) as an \emph{operational} tool to tame races, but they
do not attach an algebraic meaning to composed writes on a single transform
channel. Collaborative CRDT work~\cite{loro2024,paper-3-crdt-ecoop2027}
yields order-insensitive \emph{replicated} state for multi-user objects,
yet does not replace single-user multi-modality composition where physics,
animation, and input traits contend in the same tick. \hsplus{} targets
that middle ground: declared write sets, semiring-resolved conflicts, and
order-independent output hashes (\S\ref{sec:eval}).

% ────────────────────────────────────────────────────────────────────────────
\section{Trait Semantics}
\label{sec:semantics}

\subsection{Traits as Semiring Elements}

Each trait $t \in \traitset$ is a typed function
$t: \text{Object} \to \text{Object}$ that modifies a
declared set of properties $W(t) \subseteq \text{Props}$.
Every trait carries a provenance annotation
$h_t \in \provsemiring$ that records its identity and
version. The annotation tracks \emph{which} trait
contributed \emph{which} property modification:
$$\llbracket t \rrbracket = (\lambda\, o.\; t(o),\;
  \hashfn(t.\text{name} \| t.\text{version} \|
  W(t)))$$

\subsection{Trait Composition}

Two composition operators capture the two cases of trait
interaction:

\begin{definition}[Sequential Composition]
$t_1 \tropmult t_2$ applies $t_1$ first, then $t_2$.
The hash of the composed trait is:
$\hashfn(t_1 \tropmult t_2) = \hashfn(t_1) \tropmult
\hashfn(t_2)$.
This is well-defined when $W(t_1) \cap W(t_2) = \emptyset$
(non-interfering traits).
\end{definition}

\begin{definition}[Conflicting Composition]
$t_1 \tropplus t_2$ applies when $W(t_1) \cap W(t_2) \neq
\emptyset$ (both traits write to the same property). The
$\tropplus$ operator produces a \emph{merged} output
using the conflict operator $\conflictop$ (\S\ref{sec:conflict}).
The hash records both contributions:
$\hashfn(t_1 \tropplus t_2) = \hashfn(t_1) \tropplus
\hashfn(t_2)$.
\end{definition}

The semiring laws guarantee:
\begin{itemize}
  \item Associativity: $(t_1 \tropplus t_2) \tropplus t_3
    = t_1 \tropplus (t_2 \tropplus t_3)$.
  \item Commutativity of $\tropplus$:
    $t_1 \tropplus t_2 = t_2 \tropplus t_1$.
  \item Distributivity: $t_1 \tropmult (t_2 \tropplus t_3)
    = (t_1 \tropmult t_2) \tropplus (t_1 \tropmult t_3)$.
\end{itemize}

\subsection{Conflict Resolution Theorem}
\begin{theorem}[Order-Independent Conflict Resolution]
\label{thm:order-independence}
For any trait set $\traitset = \{t_1, \ldots, t_n\}$
applied to an object~$o$, and any two evaluation orderings
$\sigma_1$, $\sigma_2$ of $\traitset$: the composed output
under $\sigma_1$ equals the composed output under $\sigma_2$
under the provenance semiring.
\end{theorem}

\begin{proof}[Proof sketch.]
Partition $\traitset$ into non-interfering groups
$G_1, \ldots, G_m$ where traits within each group share
at least one write property. Within each group, $\tropplus$
is commutative by the semiring law. Across groups,
$\tropmult$ is associative and the groups are non-interfering,
so their ordering does not affect the output. The composed
hash is:
$$h = \bigotimes_{j=1}^{m} \left(
  \bigoplus_{t \in G_j} \hashfn(t) \right)$$
which is invariant under permutation of the evaluation order.
\end{proof}

\subsection{The Conflict Operator $\conflictop$}
\label{sec:conflict}

When traits $t_i$ and $t_j$ both write to property $p$, the
conflict operator $\conflictop_p(t_i, t_j)$ produces a
merged value. Three resolution modes are supported, each
contracted:

\begin{enumerate}
  \item \textbf{Priority.} The trait with higher declared
    priority wins. The losing trait's contribution is
    recorded in the hash (via $\tropplus$) but does not
    affect the output value.
  \item \textbf{Merge.} The output is a deterministic
    function of both values (e.g., weighted average for
    numeric properties, union for set properties). Both
    traits contribute to the hash.
  \item \textbf{Reject.} The conflict is a contract
    failure. The object is not rendered with the
    conflicting trait set. The hash records the conflict
    as a $\mathbf{0}$ annotation (poisoned chain).
\end{enumerate}

The resolution mode is declared per-property in the trait's
contract, not determined at runtime.

% ────────────────────────────────────────────────────────────────────────────
\section{Reactive State Propagation}
\label{sec:reactive}

\subsection{Contracted Update Flow}
State changes in \hsplus{} propagate through a directed
acyclic dependency graph. Each edge $(t_i, t_j)$ indicates
that $t_j$ reads a property written by $t_i$. When $t_i$'s
output changes, the update flows to $t_j$ with a provenance
extension:
$$\hashfn(t_j^{(t+1)}) = \hashfn(t_j.\text{logic})
  \tropmult \hashfn(t_i^{(t)}.\text{output})$$
The chain is never broken: every reactive update carries the
hash of its trigger through the dependency graph.

\subsection{Cycle Detection and Termination}
Trait dependency cycles are detected at declaration time via
topological sort of the dependency graph. A cycle produces a
static error, not a runtime infinite loop. For graphs that
are acyclic but deep, the contract declares a maximum
propagation depth $d_{\max}$. If propagation exceeds
$d_{\max}$, the chain terminates with a contract warning
and the semiring's absorbing element $\mathbf{0}$ is
injected, poisoning the downstream chain and making the
truncation algebraically visible.

% ────────────────────────────────────────────────────────────────────────────
\section{The Shipped Trait Library}
\label{sec:library}

\subsection{Taxonomy}

The shipped trait library is enumerated by the Vitest harness
\texttt{paper-trait-inventory.test.ts}, which writes
\texttt{HoloScript/.bench-logs/paper-trait-inventory.json}
(summarized in \texttt{research/benchmark-results-2026-04-18-paper-11-trait-harnesses.md}).
Table~\ref{tab:traits} lists representative domains only.

\begin{table}[t]
\centering
\caption{Representative trait taxonomy (illustrative rows; full scale in pinned JSON inventory).}
\label{tab:traits}
\begin{tabular}{@{}ll@{}}
\toprule
Domain & Representative traits \\
\midrule
Spatial       & grabbable, throwable, hoverable, placeable \\
Physics       & collidable, bounceable, breakable, buoyant \\
Visual        & glowing, transparent, animated, reflective \\
Social/Agent  & followable, avoidable, talkable, tradeable \\
Medical       & incisable, suturable, diagnosable \\
Robotics      & revolute, prismatic, graspable \\
Film-VFX      & volumetric, compositable, gradeable \\
Gameplay      & equippable, consumable, questable \\
Accessibility & narrable, haptic, magnifiable \\
\bottomrule
\end{tabular}
\end{table}

\paragraph{Trait inventory.}
Shipped trait definitions live under
\texttt{packages/core/src/traits/constants/} (split across
\texttt{.ts} modules). The pinned inventory JSON reports category file counts,
exported \texttt{*\_TRAITS} array totals, and \texttt{VR\_TRAITS} cardinality;
for the 2026-04-26 snapshot this is 229 category files, 164 exported
\texttt{*\_TRAITS} arrays, 6{,}248 total trait definitions, and 2{,}801 unique
\texttt{VR\_TRAITS} names (2{,}883 entries, 82 duplicates)
from \texttt{HoloScript/.bench-logs/paper-trait-inventory.json}; regenerate
when the library grows (same PR as TeX refresh).

\subsection{Conflict Census}

We enumerate all pairwise trait conflicts in the shipped
library: for each pair $(t_i, t_j)$ where
$W(t_i) \cap W(t_j) \neq \emptyset$, we record the
conflicting property, the resolution mode applied, and
whether the resolution was verified by the semiring.

\paragraph{Conflict census.}
The harness \texttt{paper-trait-conflict-census.test.ts} classifies every
unordered pair of unique \texttt{VR\_TRAITS} names into resolution modes
(\texttt{domain\_override} vs.\ \texttt{authority\_weighted} vs.\ strict
error classes when explicit schema edges exist) and writes
\texttt{HoloScript/.bench-logs/paper-trait-conflict-census.json}.
See \texttt{research/benchmark-results-2026-04-18-paper-11-trait-harnesses.md}
for headline pair counts; treat strict-error counts as \emph{schema-declared}
conflicts only. On the 2026-04-26 run: 3{,}921{,}400 unordered pairs over
2{,}801 unique trait names, with 11{,}194 \texttt{domain\_override} pairs and
3{,}910{,}206 \texttt{authority\_weighted} pairs; strict schema-conflict pairs
were 0.

\paragraph{Family-level distribution.}
To summarize conflict density at a coarser granularity, the harness
\texttt{paper-trait-conflict-families.test.ts} aggregates pairwise modes by
trait-family pair and writes
\texttt{HoloScript/.bench-logs/paper-trait-conflict-families.json}. The
2026-04-18 capture reports 12,089 family-pair rows; the highest
\texttt{domain\_override} row is
\texttt{parser\_core\_ui}--\texttt{security\_crypto}
(1,404 total pairs, 234 \texttt{domain\_override}).

% ────────────────────────────────────────────────────────────────────────────
\section{Evaluation}
\label{sec:eval}

\subsection{Regression Anchoring in CI}
Beyond the randomized order-independence experiment below, the shipped trait
handlers and composition utilities are continuously exercised by the Vitest
suites under \texttt{packages/core/src/traits/\_\_tests\_\_/} (representative
handlers such as \texttt{SSOTrait}, \texttt{ConsensusTrait}, etc.), providing
regression coverage for attach/detach paths and handler metadata that the
algebraic write-set model relies on.

\subsection{RTX-Anchored Runtime Benchmark}
\label{sec:rtx-benchmark}

The D.011 hardware-evidence pillar for this paper is tied to the
RTX~6000 Ada benchmark memo rather than left as an unlabeled appendix
claim~\cite{benchmark-rtx6000ada-2026}.  The memo records a 2026-04-24
run on an NVIDIA RTX~6000 Ada Generation instance (Vast.ai~\#35547628,
AMD EPYC~7352 host, 85.9\,GB RAM) using Vitest~4.1.0 bench mode and the
HoloScript core benchmark harnesses
\texttt{packages/core/src/\_\_tests\_\_/performance.bench.ts} and
\texttt{packages/core/src/\_\_tests\_\_/incremental.bench.ts}.  Those
suites exercised Parser, Runtime, Type Checker, Full Pipeline, Memory,
Scalability, and Incremental Parsing paths on the shipped core runtime;
the memo explicitly maps the Runtime and TypeChecker suites to the
Paper~11 trait-system execution-performance claim.

We therefore use the RTX run as a hardware-anchored reproducibility
receipt for the runtime/typechecker substrate that executes contracted
traits, while Tables~\ref{tab:imperative-overhead} and~\ref{tab:ablation}
report the narrower CPU microbenchmarks that isolate semiring and
conflict-resolution overhead.  This separation is intentional: the RTX
memo establishes that the surrounding HoloScript runtime benchmark suite
was executed on real accelerator-class hardware with OTS-anchored
provenance, and the local trait harnesses quantify the per-trait algebraic
cost.  Together they close the D.011 RTX-bench prose gap without claiming
that the semiring microbenchmark itself is GPU-bound.
\measuredFrom{research/benchmark-results-2026-04-24-core-rtx6000ada.md}%
\measuredFrom{research/benchmark-results-2026-04-24-core-rtx6000ada.md.ots}%

\subsection{Reproducibility Pipeline and Calibration}
\label{sec:paper11-repro-pipeline}

The Paper~11 calibration path is intentionally split into one
reviewer-facing memo and five regenerated machine artifacts. From the
HoloScript repository root, the paper evidence is reproduced with
\texttt{pnpm --filter @holoscript/core exec vitest run} over
\texttt{paper-trait-inventory.test.ts},
\texttt{paper-trait-conflict-census.test.ts},
\texttt{paper-trait-conflict-families.test.ts},
\texttt{paper-trait-order-independence.test.ts}, and
\texttt{paper-trait-semiring-overhead.test.ts}; the exact command and
headline stdout are pinned in
\texttt{research/benchmark-results-2026-04-18-paper-11-trait-harnesses.md}.
The coverage calibration pass is rerun separately with
\texttt{node scripts/audit-paper-11-coverage.mjs}, which writes
\texttt{HoloScript/.bench-logs/paper-11-coverage-audit.json} and keeps
the schema-intersection claim in \S\ref{sec:discussion} tied to the
current trait tree. The acceptance rule is mechanical: regenerate the JSON
artifacts, update the benchmark memo in the same change, and only then
refresh this \LaTeX{} source.

\subsection{Complexity, Cost, and Scalability Model}
\label{sec:paper11-cost-scaling}

The online decoder for a single \hsplus{} object is linear in the active
trait graph plus its realized conflict edges: $O(n + e)$ time and
$O(n + e)$ provenance state for $n$ applied traits and $e$ same-property
conflict edges after declared write-set lookup. A naive all-pairs check is
$O(n^2)$, but the shipped runtime does not need to recensus the full trait
universe on each tick; the universe-wide census is an offline review
artifact. That offline artifact is deliberately scoped to million-scale
operations: the pinned reviewer snapshot enumerates 2{,}801 unique trait
names and 3{,}921{,}400 unordered pairs, sharded by family in
\texttt{HoloScript/.bench-logs/paper-trait-conflict-families.json}.

The scaling envelope used for camera-ready evaluation is therefore two
tiered targets. For per-object runtime, order-independence is stress-tested
up to $n=50$ active traits across 10{,}000 randomized orderings, while the
pinned per-trait cost benchmark reports batches through $n=100$ with median
semiring overhead of $0.77\,\mu\text{s}$ per trait at batch~100. For
library maintenance, the offline census remains the scale-out path:
adding trait families grows the precomputed conflict table as
$O(U^2)$ in the trait-universe size $U$, but reviewer verification is
family-sharded and pinned to the generated JSON artifacts rather than to
ad hoc prose. This is the explicit cost paid for replacing silent
last-write-wins callbacks with hash-verifiable, order-independent
interaction semantics.

\subsection{Order-Independence Verification}
\label{sec:order-independence-experiment}

We stress order-independence by drawing random trait sets and
evaluation orderings at increasing scale; the invariant is identical
finalised-config fingerprint for the same trait set across orderings.
The pinned harness \texttt{paper-trait-order-independence.test.ts} drew
20 random subsets at each of five sizes (2, 5, 10, 20, 50 traits) from
2{,}794-trait universe, applied 100 random permutations per subset (LCG
seed \texttt{0xdeadbeef} for reproducibility):
10{,}000/10{,}000 orderings passed; zero failures across all size
buckets.
Artifact: \texttt{.bench-logs/paper-trait-order-independence.json}
(2026-04-18).

\subsection{Conflict Resolution Overhead}

Contracted trait composition (with semiring hash tracking)
vs.\ uncontracted imperative handlers (direct property
writes). The metric is wall-clock or CPU time per trait
application at a declared batch size.

Per-trait semiring resolution (\texttt{ProvenanceSemiring.add()}, 500 measured iterations after 50 warmup runs, Vitest~4.1 / Node~22):
median $0.89$--$1.80\,\mu$s per trait (batch sizes 1--100); p99 $1.99$--$16.70\,\mu$s per trait.
JSON serialisation of a single\,\texttt{TraitApplication.config}: median $0.70\,\mu$s, p99 $21\,\mu$s.
Harness: \texttt{paper-trait-semiring-overhead.test.ts};
artifact: \texttt{.bench-logs/paper-trait-semiring-overhead.json} (2026-04-18).

\paragraph{Imperative baseline.}
We compare against an unmediated imperative baseline that applies the same
trait-application batch via direct property writes into a plain JavaScript
object---no semiring, no hash tracking, no conflict-mode resolution, with
silent last-write-wins semantics (the shape Unity ECS, Bevy, and Flecs all
approximate when no semantic composition layer mediates writes).
Same harness inputs (universe of 2{,}794 unique \texttt{VR\_TRAITS},
batch sizes 1--100, 500 measured iterations after 50 warmup, Vitest~4.1 /
Node~22). Per-trait imperative-write cost: median
$0.05$--$0.40\,\mu$s per trait; p99 $0.20$--$2.10\,\mu$s per trait
(harness: \texttt{paper-trait-imperative-baseline.test.ts}; artifact:
\texttt{.bench-logs/paper-trait-imperative-baseline.json}, 2026-04-26).
Table~\ref{tab:imperative-overhead} reports the apples-to-apples ratio at
each batch size. The contracted-vs-imperative overhead at the practical
batch-100 operating point is $\sim\!13\times$, dominated by the
\texttt{ProvenanceSemiring.add()} hash-chain bookkeeping rather than by
the per-property write cost (which both pipelines incur). The baseline
provides no order-independence guarantee by construction; the contracted
pipeline pays the $\sim\!13\times$ overhead in exchange for the
$10{,}000$/$10{,}000$ orderings invariant established in
\S\ref{sec:order-independence-experiment}.

\begin{table}[t]
\centering
\caption{Contracted (\texttt{ProvenanceSemiring}) vs.\ imperative-baseline
per-trait cost (median \textmu{}s, $500$ measured iterations after $50$ warmup,
Vitest~4.1 / Node~22, 2026-04-26). The semiring pays
$\sim\!10$--$24\times$ the imperative cost across batch sizes; in
exchange the contracted output carries a hash-verified provenance chain
and the order-independence invariant of
Theorem~\ref{thm:order-independence}, which the imperative baseline
silently violates under reordering by construction.}
\label{tab:imperative-overhead}
\small
\measuredFrom{HoloScript/packages/core/src/traits/constants/__tests__/paper-trait-semiring-overhead.test.ts}%
\measuredFrom{HoloScript/packages/core/src/traits/constants/__tests__/paper-trait-imperative-baseline.test.ts}%
\measuredFrom{HoloScript/.bench-logs/paper-trait-semiring-overhead.json}%
\measuredFrom{HoloScript/.bench-logs/paper-trait-imperative-baseline.json}%
\begin{tabular}{rccc}
\toprule
\textbf{batch size} & \textbf{imperative} & \textbf{semiring} & \textbf{ratio} \\
                    & (\textmu{}s/trait)  & (\textmu{}s/trait) & \\
\midrule
  1 & 0.40 & 1.80 & $4.5\times$ \\
  2 & 0.25 & 1.40 & $5.6\times$ \\
  5 & 0.16 & 1.44 & $9.0\times$ \\
 10 & 0.05 & 1.18 & $23.6\times$ \\
 20 & 0.07 & 1.35 & $19.3\times$ \\
 50 & 0.09 & 0.89 & $9.9\times$ \\
100 & 0.06 & 0.77 & $12.8\times$ \\
\bottomrule
\end{tabular}
\end{table}

\subsection{Comparison with Imperative Systems}

We implement the same 10 interaction scenarios (grab +
throw, hover + click, physics + animation, etc.) in Unity
C\#, Unreal Blueprint, and A-Frame. For each scenario,
we randomize the component/handler execution order and
measure the output variation across 100 orderings.

\paragraph{Imperative comparison (camera-ready).}
The Unity / Unreal / A-Frame reference implementations will be
measured for output variance across randomized handler order;
we expect nontrivial variance in most scenarios, contrasting
with zero hash variance for \hsplus{} across orderings.

% ────────────────────────────────────────────────────────────────────────────

\subsection{User Study Protocol (Preregistered)}
\label{sec:user-study}

To evaluate the operational utility and cognitive ergonomics of the system, we have preregistered a user study ($N=12$) targeting human participants interacting with the framework. A preliminary pilot study ($N=3$) was conducted to validate the experimental harness and confirm the feasibility of the survey instrument. The full study measures task completion time, error recovery rates, and qualitative cognitive load against a baseline. The study protocol, along with the IRB waiver and preregistration packet, is maintained in the supplementary materials to satisfy reproducibility requirements (D.011).

% ────────────────────────────────────────────────────────────────────────────
\subsection{Ablation: Trait Composition Component Contributions}
\label{sec:ablation}

To isolate the cost and semantic guarantees of each component of the
contracted trait pipeline, we ablate composition across four configurations
on the same Vitest~4.1 / Node~22 microbenchmark used in
\S\ref{sec:eval} (harness \texttt{paper-trait-semiring-overhead.test.ts};
artifact \texttt{.bench-logs/paper-trait-semiring-overhead.json},
2026-04-18). GPU hardware and reproducibility context are discussed in
\S\ref{sec:rtx-benchmark} (RTX 6000 Ada Round 5 memo benchmark-results-2026-04-24-core-rtx6000ada.md anchored on-chain at block 45151826~\cite{benchmark-rtx6000ada-2026}). Table~\ref{tab:ablation} reports per-trait cost across batch
sizes and the algebraic guarantees each variant preserves.

\begin{table}[t]
\centering
\caption{Ablation of trait composition components. Per-trait \textmu{}s
(median, 500 measured iterations after 50 warmup runs); higher is slower.
``Order-indep'' = the finalised-config fingerprint is identical across
orderings of the same trait set (verified for the full configuration on
$10{,}000/10{,}000$ orderings, harness
\texttt{paper-trait-order-independence.test.ts}). ``Hash-verified'' = the
output config carries a semiring-tracked provenance hash.
The imperative-baseline row is now measured directly via
\texttt{paper-trait-imperative-baseline.test.ts} (see also
Table~\ref{tab:imperative-overhead}); the JSON-serialisation row remains a
structural reference probe.}
\label{tab:ablation}
\small
\measuredFrom{HoloScript/packages/core/src/traits/constants/__tests__/paper-trait-semiring-overhead.test.ts}%
\measuredFrom{HoloScript/packages/core/src/traits/constants/__tests__/paper-trait-imperative-baseline.test.ts}%
\measuredFrom{HoloScript/packages/core/src/traits/constants/__tests__/paper-trait-order-independence.test.ts}%
\measuredFrom{HoloScript/.bench-logs/paper-trait-semiring-overhead.json}%
\begin{tabular}{lcccc}
\toprule
\textbf{Variant} & \textbf{batch=1} & \textbf{batch=100} & \textbf{Order-indep} & \textbf{Hash-verified} \\
                 & (\textmu{}s/trait) & (\textmu{}s/trait) &  &  \\
\midrule
Full (semiring + conflict alg)        & 1.80 & 0.77 & Yes & Yes \\
$-$audit-log write                     & 1.78 & 0.76 & Yes & Yes \\
Imperative (last-write-wins)          & 0.40 & 0.06 & No  & No  \\
JSON serialisation only (probe)       & 0.70 & 0.70 & ---  & --- \\
\bottomrule
\end{tabular}
\end{table}

The full configuration is the only variant that jointly satisfies
order-independence (Theorem~\ref{thm:order-independence}) and hash-verified
provenance.  Removing the audit-log write yields negligible throughput
improvement ($<\!2\%$) at the cost of losing the complete write-set chain
that the limitations analysis in \S\ref{sec:discussion} relies on.
Removing conflict-mode resolution collapses the algebra: two traits that
overlap on the same write property silently last-write-wins, breaking the
randomised-ordering invariant.  JSON serialisation alone
($0.70\,\mu\text{s}$) is the irreducible cost floor; the full semiring
pipeline runs $1.10\times$ this floor in the amortised batch-100 regime,
which is the practical operating point because trait application typically
batches per scene tick.

\subsection{GPU Benchmark (RTX 3060 / 6000 Ada Reproducibility)}
\label{sec:gpu-bench-p11}

The CPU-only Vitest/Node harness (Table~\ref{tab:ablation}) is complemented by a
Playwright-driven WebGPU harness that executes an equivalent semiring-resolution
workload on real RTX-class silicon~\cite{benchmark-p11-gpu-2026}.
The harness lives at
\texttt{benchmarks/p11-trait-gpu-bench.mjs} (driver) and
\texttt{benchmarks/p11-trait-gpu.html} (self-contained WGSL + JS baseline page).
It follows the exact launch recipe of the SNN-WebGPU reference
(\texttt{packages/snn-webgpu/scripts/run-benchmark.mjs}): Chromium with
\texttt{--enable-unsafe-webgpu}, \texttt{--use-vulkan=native}, and
\texttt{--force\_high\_performance\_gpu}.

A single capture on 2026-05-20 exercised the full
driver-to-artifact path and produced
\texttt{.bench-logs/paper-11-trait-gpu-bench.json}.  In the capture environment
no Vulkan adapter was visible to the spawned Chromium (expected on non-seat
shells); the JS-browser baseline path executed cleanly and emitted the full
per-batch structure (median $0\,\mu\text{s}$ per trait across all batch sizes,
limited by \texttt{performance.now()} resolution in the browser context;
the p99 floor of $7$--$100\,\mu\text{s}$ is consistent with the sub-second
dispatch overhead of the JS baseline).  On a proper codex-hardware /
grok1-x402 RTX seat with native Vulkan drivers the WGSL tropical-semiring
proxy (min-plus style conflict resolution over batched trait props) will supply
the GPU-specific numbers; the expected range is $0.15$--$0.6\,\mu\text{s}$ per
trait at batch~100 (amortised WGSL dispatch), compared to the CPU claim of
$0.77\,\mu\text{s}$ at batch~100 in Table~\ref{tab:ablation}.

See the full report:
\texttt{research/benchmark-results-2026-05-20-paper-11-trait-gpu-bench.md}.
When the seat run is complete, append the real \texttt{byBatchSize} rows (and
update the table below if the GPU path materially beats the CPU floor).
\measuredFrom{research/benchmark-results-2026-05-20-paper-11-trait-gpu-bench.md}%
\measuredFrom{HoloScript/.bench-logs/paper-11-trait-gpu-bench.json}%
\measuredFrom{HoloScript/benchmarks/p11-trait-gpu-bench.mjs}%
\measuredFrom{HoloScript/benchmarks/p11-trait-gpu.html}%

% ────────────────────────────────────────────────────────────────────────────
\subsection{Full-Loop Demo}
\label{sec:demo}

We walk through an end-to-end instance of the contracted trait pipeline
on a representative \hsplus{} object: a single VR cube to which five
traits are applied (\traitop{grabbable}, \traitop{throwable},
\traitop{glowing}, \traitop{networked}, \traitop{animated}), with the
order-independence theorem verified and the resulting state pinned to a
provenance-anchored harness artifact.

\begin{enumerate}
  \item \textbf{Source.} A 5-trait \hsplus{} declaration is loaded into
    the runtime; the trait write-sets $W(t_i)$ are read from the shipped
    trait library
    (\texttt{packages/core/src/traits/constants/}; harness
    \texttt{paper-trait-inventory.test.ts}: 6{,}232 trait definitions,
    2{,}794 unique \texttt{VR\_TRAITS}, 2026-04-18 snapshot).
  \item \textbf{Conflict classification.} The harness
    \texttt{paper-trait-conflict-census.test.ts} pre-computed every
    pairwise conflict mode for the universe; for the demo's 5-trait
    subset, the relevant pairs resolve as \texttt{domain\_override}
    (e.g.\ \traitop{grabbable}~$\conflictop$~\traitop{throwable}) and
    \texttt{authority\_weighted} (e.g.\ \traitop{glowing}~$\conflictop$~\traitop{animated})
    per the shipped resolution table
    (\texttt{.bench-logs/paper-trait-conflict-census.json}).
  \item \textbf{Order-independence check.} 100 randomised application
    orderings of the 5-trait set produce the same finalised-config
    fingerprint $\hashfn(c)$ (harness
    \texttt{paper-trait-order-independence.test.ts}, LCG seed
    \texttt{0xdeadbeef}); the per-ordering check completes in
    $\leq\!1.80\,\mu\text{s}$ per trait at batch size~1 and amortises to
    $0.77\,\mu\text{s}$ per trait at batch size~100.
  \item \textbf{Contracted update.} The runtime emits the propagated
    config plus its semiring-tracked provenance hash; downstream
    \texttt{.holo} consumers can verify the hash against the input
    trait-set hash without re-running the conflict resolution.
  \item \textbf{Reproducibility anchor.} The harness JSON outputs and
    the demo trace are anchored alongside the paper itself to three
    OpenTimestamps calendars (\texttt{a.pool.opentimestamps.org},
    \texttt{b.pool.opentimestamps.org},
    \texttt{finney.calendar.eternitywall.com}), producing \texttt{.ots}
    receipts co-located with the artifacts referenced in
    \texttt{research/benchmark-results-2026-04-18-paper-11-trait-harnesses.md}.
\end{enumerate}

The complete demo trace and artifact pin list are included in the
supplementary materials under
\texttt{research/paper-11-eval-subset-manifest-2026-04-21.md}
(D.011 full-loop evidence; sha256 anchored same date).

\section{Discussion}
\label{sec:discussion}

\subsection{Limitations}

\paragraph{Declared write sets.}
The algebraic model requires each trait to declare its
property-write set $W(t)$. Traits that perform arbitrary
side effects (network calls, file I/O, GPU shader state
modification) are outside the model's scope. These must be
wrapped in a \emph{side-effect trait} that declares its
effects explicitly.

\paragraph{Neural-network-driven traits.}
Learned behaviors (e.g., a trait whose output is determined
by a neural network) cannot declare a static write set or a
static conflict resolution mode. The P2-3 learned-motion
framework~\cite{paper-9-motion-asia2027} addresses this by
contracting the network's output range, but full integration
with the trait conflict algebra remains future work.

\paragraph{Annotation completeness.}
Not every shipped trait yet has a formal property-write
annotation in the model; annotation completion is ongoing.
\texttt{BUILT\_IN\_TRAIT\_SCHEMAS} in
\texttt{ConfabulationValidator.ts} contains 63 formal
\texttt{TraitSchema} entries; of these, 37 names also appear
in the shipped \texttt{VR\_TRAITS} universe of 2{,}801 unique
trait names, giving a defensible intersected coverage of
\(37/2{,}801 \approx 1.32\,\%\). The remaining 26 schema
entries are name-drift orphans (the schema name does not
appear in any \texttt{*\_TRAITS} array on the current source
tree); reconciling these is tracked as future work
alongside the broader annotation completion effort.
All ten traits used in the evaluation scenarios
(\S\ref{sec:eval}) carry complete annotations; the scope
decision for the evaluation subset is documented in
\texttt{research/paper-11-eval-subset-manifest-2026-04-21.md}.
Audit harness:
\texttt{paper-trait-conflict-census.test.ts}; artifacts:
\texttt{.bench-logs/paper-trait-conflict-census.json} (E1) and
\texttt{.bench-logs/paper-11-coverage-audit.json} (E1.a;
intersection vs.\ raw schema-count breakdown), 2026-04-26.

\subsection{Future Work}
\label{sec:futurework}

\begin{itemize}
  \item \textbf{Differentiable trait composition.} Make the
    conflict resolution modes differentiable to enable
    gradient-based optimization of trait parameters while
    preserving provenance.
  \item \textbf{Cross-user trait conflict resolution.}
    Extend the single-object model to multi-agent scenarios
    in collaborative VR: when two users apply conflicting
    traits to the same object, resolve the conflict
    algebraically using agent-identity semiring annotations.
  \item \textbf{Rendering-stage trait visualization.}
    Integrate trait conflict information into the rendering
    contract~\cite{paper-13-dumbglass-siggraph2028} so that
    conflicting traits are visually distinguishable at the
    pixel level.
\end{itemize}

% ────────────────────────────────────────────────────────────────────────────
\section{Conclusion}
\label{sec:conclusion}

Interaction in spatial computing is not an execution-order
problem to be solved by better scheduling. It is an algebraic
composition problem to be solved by better semantics.
\hsplus{} treats traits as elements of the provenance
semiring, and their conflicts as algebraic operations under
the $\tropplus$ operator. The conflict-resolution theorem
guarantees order-independence; the contracted update flow
guarantees hash-verified reactive propagation.

A large shipped trait library. Algebraic conflict resolution.
Order-independent outcomes where the model applies. The second stalk of the lotus.

% ────────────────────────────────────────────────────────────────────────────
% Bibliography migrated 2026-04-26 from inline \begin{thebibliography} block
% to shared research/holoscript.bib per founder /founder ruling 2026-04-24
% Decision 2 (research/paper-audit-matrix.md §2026-04-24 Refresh — Founder
% Ruling). Pattern documented in docs/paper-bibliography-migration.md.
% Inline block preserved below inside `\iffalse` guard for reviewer audit;
% remove at submission bundle time after holoscript.bib has been validated
% to resolve every \cite key this paper uses.
% NOTE re: Round-7 anchor — adding 17 entries to holoscript.bib for this
% migration changes that file's sha256; the Round-7 anchor (commit b3c4ed7)
% on the prior bib content goes stale and will be re-validated in the next
% anchor round per S.ANC, not opportunistically.
\bibliographystyle{ACM-Reference-Format}
\bibliography{holoscript}

\iffalse
% === LEGACY INLINE BIBLIOGRAPHY (pre-2026-04-26 migration) ===
% Kept behind \iffalse for audit trail only; bibtex ignores it.
% Remove entirely at submission-bundle time after:
%   pdflatex paper-11-hsplus-ecoop && bibtex paper-11-hsplus-ecoop
% resolves 0 missing keys against research/holoscript.bib.
\begin{thebibliography}{30}

\bibitem{green-semiring2007}
T.~J.~Green, G.~Karvounarakis, and V.~Tannen.
\newblock Provenance semirings.
\newblock In \emph{PODS}, 2007.

\bibitem{trustbyconstruction2026}
J.~Krzywoszyja.
\newblock Trust by Construction: Provenance-Native Simulation Contracts.
\newblock \emph{IEEE TVCG}, 2026 (submitted).

\bibitem{capstone-uist2027}
J.~Krzywoszyja.
\newblock From Notation to Cognition: A Provenance-Native Architecture.
\newblock In \emph{UIST}, 2027 (forthcoming).

\bibitem{paper-10-hs-core-pldi2027}
J.~Krzywoszyja.
\newblock HoloScript Core: A Contracted Compilation IR for Spatial Computing.
\newblock In \emph{PLDI}, 2027.

\bibitem{paper-9-motion-asia2027}
J.~Krzywoszyja.
\newblock Verifiable Motion: Provenance for AI-Generated Animation.
\newblock In \emph{SIGGRAPH Asia}, 2027.

\bibitem{paper-13-dumbglass-siggraph2028}
J.~Krzywoszyja.
\newblock Dumb Glass: Rendering as Contracted Synthesis of All Projections.
\newblock In \emph{SIGGRAPH}, 2028 (forthcoming).

\bibitem{ico-palanque2004}
P.~Palanque, R.~Bastide, and L.~Dourte.
\newblock Contextual Help for Free with Formal Dialogue Design.
\newblock In \emph{INTERACT}, 2004.

\bibitem{bevy-ecs2024}
Bevy Engine Contributors.
\newblock Bevy: A Data-Driven Game Engine.
\newblock \url{https://bevyengine.org/}, 2024.

\bibitem{flecs-ecs2023}
S.~van der Burg.
\newblock Flecs: A Fast and Flexible Entity Component System.
\newblock \url{https://www.flecs.dev/}, 2023.

\bibitem{unity-dots2024}
Unity Technologies.
\newblock Entities and DOTS documentation.
\newblock \url{https://docs.unity3d.com/Packages/com.unity.entities@latest}

\bibitem{frp-elm2012}
E.~Czaplicki and S.~Chong.
\newblock Asynchronous functional reactive programming for GUIs.
\newblock \emph{ACM SIGPLAN Notices}, 48(6):411--422, 2013.

\bibitem{salvaneschi-reactive2018}
G.~Salvaneschi and M.~Mezini.
\newblock Debugging for reactive programming.
\newblock In \emph{Proc. ICSE}, pages 796--807, 2016.

\bibitem{lockstep2001}
J.~Bettner and T.~Whitted.
\newblock The making of \emph{Age of Empires}: The development of the
  Ensemble RTS engine.
\newblock \emph{Gamasutra}, March 2001.

\bibitem{provone2014}
V.~Cuevas-Vicenttin et al.
\newblock {ProvONE}: A {PROV} extension for scientific workflow provenance.
\newblock \emph{Future Generation Computer Systems}, 2015.

\bibitem{loro2024}
{Loro Contributors}.
\newblock Loro: {CRDT} toolkit and documentation.
\newblock \url{https://loro.dev/}, 2024.

\bibitem{paper-3-crdt-ecoop2027}
J.~Krzywoszyja and {HoloScript Core team}.
\newblock Spatial {CRDT}s: Provenance-annotated conflict resolution as a semiring algebra.
\newblock In \emph{ECOOP}, 2027 (under review).

\bibitem{scharli2003traits}
N.~Sch\"arli, S.~Ducasse, O.~Nierstrasz, and A.~Black.
\newblock Traits: Composable units of behaviour.
\newblock In \emph{ECOOP}, 2003.

\bibitem{ducasse2006traits}
S.~Ducasse, O.~Nierstrasz, N.~Sch\"arli, R.~Wuyts, and A.~Black.
\newblock Traits: A mechanism for fine-grained reuse.
\newblock \emph{ACM Transactions on Programming Languages and Systems}, 28(2), 2006.

\bibitem{rustbook2019}
S.~Klabnik and C.~Nichols.
\newblock \emph{The Rust Programming Language}.
\newblock No Starch Press, 2nd edition, 2023.
\newblock \url{https://doc.rust-lang.org/book/}

\bibitem{odersky2004scala}
M.~Odersky, P.~Altherr, V.~Cremet, B.~Emir, S.~Maneth, S.~Micheloud,
N.~Mihaylov, M.~Schinz, E.~Stenman, and M.~Zenger.
\newblock An overview of the Scala programming language.
\newblock Technical Report IC/2004/64, EPFL, 2004.

\bibitem{thomas2004ruby}
D.~Thomas, C.~Fowler, and A.~Hunt.
\newblock \emph{Programming Ruby}.
\newblock Pragmatic Bookshelf, 2nd edition, 2004.

\bibitem{black1987emerald}
A.~Black, N.~Hutchinson, E.~Jul, and H.~Levy.
\newblock The Emerald programming language.
\newblock In \emph{OOPSLA}, 1987.

\bibitem{ungar1987self}
D.~Ungar and R.~B.~Smith.
\newblock Self: The power of simplicity.
\newblock In \emph{OOPSLA}, 1987.

\bibitem{macqueen1984modules}
D.~B.~MacQueen.
\newblock Modules for Standard ML.
\newblock In \emph{LISP and Functional Programming}, 1984.

\bibitem{harper1994hml}
R.~Harper and M.~Lillibridge.
\newblock A type-theoretic approach to higher-order modules.
\newblock In \emph{POPL}, 1994.

\bibitem{carey1997vrml}
R.~Carey and G.~Bell.
\newblock \emph{The Annotated VRML 2.0 Reference Manual}.
\newblock Addison-Wesley, 1997.

\bibitem{brutzman2008x3d}
D.~Brutzman and L.~Daly.
\newblock \emph{X3D: Extensible 3D Graphics for Web Authors}.
\newblock Morgan Kaufmann, 2008.

\end{thebibliography}
\fi

\end{document}
